\chapter{El juicio del comandante Rykov}\label{sec:el-juicio-del-comandante-rykov}

Mientras me acompañaban por el Salón de Gobierno, aminoré el paso para admirar las obras de arte distribuidas por el recinto: murales que representaban el cosmos, esculturas de deidades antiguas, artefactos culturales de hacía eones...

—Comandante Rykov —el guardia yendil pronunció mi nombre lentamente, ya que su especie tenía dificultades con el sonido de la «r»—. ¿Hay alguna razón por la que se haya detenido?

—Solo estoy mirando. Es fascinante, una muestra de todo... —negué con la cabeza.

—Puede ver las exposiciones en otro momento. El Senado lo está esperando —gruñó.

—Bien.

Tal vez debí sentirme irritado por la impaciencia de ese tipo, pero, francamente, me alegraba que no me tuviera miedo. La mayoría de los extraterrestres con los que había hablado desde mi llegada andaban con pies de plomo, aterrorizados de decir algo incorrecto. Pese a siglos de moderación y no violencia, pensaban que los humanos perderíamos los estribos por capricho.

¿Por qué el Comando Terrano había insistido en que testificara ante el Senado? ¿Por qué molestarse en cumplir con su citación? Esto no sería más que otra farsa, la última inquisición contra la humanidad.

Forcé una sonrisa mientras caminaba hacia la cámara central, sintiendo cómo todas las miradas se posaban en mí. El nuevo presidente estaba de pie ante el atril, revisando un informe. Su nombre era Retke, exembajador coviano. Su especie estaba en la mitad de la escala de agresividad, y él era uno de los representantes más jóvenes del Senado.

Aparte de eso, no sabía mucho sobre él, pero no esperaba que nos diera una oportunidad justa. Esta era su oportunidad de oro para demostrar su valía y ganar puntos políticos con los partidarios de Ula.

—Comandante Rykov, del planeta Tierra. Por favor, tome asiento en el banco designado —la voz de Retke era suave y sedosa, agradable al oído—. Estoy seguro de que comprende por qué está aquí.

—Gracias, señor presidente. Sé de qué desea hablar —respondí.

Retke dio unos golpecitos con sus garras en el atril. —Muy bien. Han sido unas semanas interesantes, ¿no cree?

—Esa es una forma de decirlo —suspiré.

—Creo que muchos de nuestros oyentes tienen preocupaciones válidas sobre los eventos ocurridos. ¿No consideran sus armas nanométricas crueles y excesivas?

—Depende de contra quién se utilicen.

 —¿Y quién decide quién merece ese destino, comandante? No creo que todos los humanos sean malvados, como creía mi predecesora... pero algunos de ustedes lo son. Ustedes lo saben. ¿Qué sucede si sus armas caen en las manos equivocadas?

Por mucho que quisiera discutir, este tipo tenía razón. Bastaba con unos cuantos idiotas de gatillo fácil en la cadena de mando para desencadenar el Armagedón; la misma razón por la que la proliferación nuclear casi causó la extinción humana. Era demasiado fácil juzgar a la ligera, como casi hicimos con los devoradores al principio de este lío, antes de conocer el panorama completo.

—Son preguntas válidas, señor presidente. Pero déjeme preguntarle algo: ¿qué pasará cuando los próximos devoradores vengan por nosotros? —hice una pausa para que mis palabras calaran—. ¿Cómo protegemos a la Federación sin algún tipo de disuasión?

—Responder a una pregunta con otra. Debería postularse para un cargo, comandante; tiene un talento natural —los representantes soltaron algunas risas—. Entiendo que es prudente tener un... último recurso, como dice. Pero son necesarios mecanismos de seguridad, supervisión y transparencia. Si realmente desean la paz, estas armas no deben estar en manos de una sola especie.

—Esa no es mi decisión, señor presidente. ¿Hay algo más que pueda responder? —entrecerré los ojos.

—Sí, sí. Su gobierno no ha ofrecido una explicación suficiente sobre varios asuntos. Empecemos por la pregunta fácil... ¿dónde está el general Kilon? Su embajador afirmó que era pasajero en su nave y desde entonces nadie ha sabido nada de él.

Kilon. Tuve la suerte de que aceptara el puesto de primer oficial en mi nave, dada su inteligencia y capacidad. Pero, aunque no se quejaba de su situación, al menos conmigo, sabía que su corazón estaba en su hogar. Era culpa mía que nunca volviera a ver su planeta. Lo mínimo que podía hacer era preservar su honor y dejar que pasara a los libros de historia como un héroe.

—Está muerto —murmullos inquietos recorrieron la cámara y le lancé una mirada de disculpa al embajador jatari, Pallum—. Recibió un disparo en el abdomen durante nuestro enfrentamiento con los xanik. No pudimos hacer nada. Por favor, acepte mis más sinceras condolencias. Era un buen soldado y un mejor hombre.

Los ojos del presidente Retke me examinaron con rapidez; parecía evaluar mi lenguaje corporal. La verdad no solo destrozaría a Kilon, sino que implicaría a la Unión Terrana en experimentos genéticos ilegales. Si había alguna esperanza de salvar nuestra relación con la Federación, esta era la historia que debían creer.

Incliné la cabeza y cerré los ojos. Las imágenes de la capital de los devoradores flotaban en mi mente, tan vívidas como cuando estuve en el puente hace una semana. Los mismos recuerdos que me asaltaban cada noche al intentar dormir. En lugar de alejarlos, me concentré en ellos.

Los océanos inmensos y los débiles resplanderes anaranjados que salpicaban las tierras gritaban que había vida. Imaginé a los niños presenciando su último amanecer, sin sospechar el destino que caería sobre su raza. La ola de calor vaporizando la superficie, extinguiendo millones de almas en un instante.

Un pueblo condenado por su propia creación. Tal vez en otra vida habríamos sido amigos.

O tal vez en otra vida habríamos sufrido su destino.

—¡Comandante Rykov! —el presidente gritó mi nombre, sacándome de mi trance. Sus fosas nasales se dilataron en un gesto de preocupación—. No quise incomodarlo. Cambiaremos de tema.

Di un silencioso suspiro de alivio. Retke necesitó ver una muestra genuina de dolor para convencerse, y esa fue la única forma en que pude dársela.

—Siguiente asunto. Me gustaría empezar expresando mi gratitud en nombre de la Federación. Si no hubiera intervenido, creo que ahora seríamos propiedad de la República Xanik —dijo.

—Fue un placer —dudé. Algo me decía que venía un «pero».

—Dicho esto... su gobierno afirma que el embajador Cazil murió por fuego amigo. Me inclinaría más a creerlo si todos los individuos en la línea de sucesión xanik no hubieran aparecido muertos al día siguiente de la incursión en la capital —murmuró Retke.

 —¿Es eso una acusación? Usted sabe tan bien como yo que hubo un golpe de Estado. Los ciudadanos comunes de Xanik tomaron el palacio y destituyeron a sus líderes. Se transmitió en vivo a toda la galaxia.

—Ojalá fuera así de sencillo. Esos eran los «ciudadanos comunes» mejor armados que he visto. ¿Los insurgentes normales tienen drones francotiradores y rifles automáticos en su mundo, comandante Rykov?

—Bueno... —la respuesta era afirmativa, pero admitirlo no pintaría bien a la humanidad—. No, supongo que no. ¿A qué se refiere?

—Digo que alguien los armó. Alguien también le dio a la población una razón para rebelarse, inundando su economía con criptomonedas y haciendo que no pudieran pagar lo básico —el presidente Retke hizo una pausa y me miró a los ojos—. Sería lógico que la misma entidad fuera responsable de ambas cosas, y la Unión Terrana ya ha sido vinculada al ataque financiero.

Bajé la mirada. —No digo que lo hicieran los humanos, pero si lo hubiéramos hecho... habría sido por el bien común. Durante siglos, la administración xanik ha sido implacable en su expansión y conquista. Su gobierno necesitaba cambios drásticos para detener esto.

Una serie de susurros ansiosos recorrieron la asamblea; mis palabras eran prácticamente una confesión. Miré a la audiencia y vi a la embajadora Johnson en su asiento, negando con la cabeza. Era evidente que no aprobaba mi respuesta, pero ¿qué más podía decir? Este nuevo presidente era demasiado astuto y ya había atado cabos. Negarlo nos haría parecer mentirosos y daría más motivos de desconfianza.

Retke esperó a que el rumor cesara antes de responder: —Es una forma de pensar ingenua y simplista.

 —¿Qué quiere decir con eso? —pregunté.

—Ustedes suponen que el próximo gobierno será mejor, pero no consideran la posibilidad de que sea peor —respondió—. El antiguo régimen fue elegido por los ciudadanos y cumplía la voluntad del pueblo. No se pueden resolver problemas sistémicos matando a unos cuantos individuos.

Me di cuenta de que estaba de acuerdo con él, a pesar de mi lealtad. El modus operandi de la inteligencia terrana me parecía absurdo; se lo había dicho a mi hermano por años. Al derrocar a la República Xanik, habían creado un vacío de poder. Muy probablemente, estaban reemplazando a un agresor democrático por uno autoritario.

Solté un suspiro cansado. —El tiempo lo dirá. ¿Eso es todo?

—Una pregunta más. ¿Qué pasó con los devoradores? —preguntó.

—Su estrella se volvió nova y su sistema fue destruido.

Técnicamente, era la verdad.

 —¿Tres mil millones de años antes de su fin natural? ¿Y justo cuando una flota se dirigía hacia la Tierra? Un momento oportuno.

—Tal vez fue intervención divina, señor presidente.

—No lo creo. Esta es su última oportunidad de confesar lo que sucedió, comandante. No me obligue a actuar.

—No sé de qué está hablando.

—Entonces se lo mostraré. Si prestan atención al proyector, verán imágenes de una sonda secreta que monitoreaba el sistema de los devoradores.

El horror recorrió mis venas al iniciarse el video detrás del presidente. Mostraba una nave solitaria aproximándose a una estrella luminosa. Aunque el exterior estaba carbonizado, el emblema terrano seguía siendo visible. No habría importado si no lo fuera; el diseño del buque insignia era inconfundible.

El Senado observó en un silencio gélido cómo la nave depositaba una cápsula en la corona estelar. La gravedad hizo el resto, atrayendo el proyectil a la superficie. Sentí náuseas y quise rogar a Retke que detuviera la grabación.

En cambio, me quedé allí sentado, inmóvil. El presidente adelantó la grabación en su holopad hasta los últimos momentos de la estrella. Los gases ardientes giraban hacia el núcleo como agua en un desagüe. Con un estremecimiento, el orbe colapsó sobre sí mismo hasta desaparecer.

Por unos segundos, solo hubo oscuridad. El universo parecía en paz.

Luego se produjo un destello blanco cegador que saturó la cámara. La imagen duró apenas un instante antes de cortarse. Sobra decir que la sonda fue destruida.

—Entonces, dígame: ¿era esa otra nave que casualmente se parecía a la suya, comandante? —el asco brilló en el rostro de Retke antes de recuperar la compostura—. Sabe, sería bueno recibir una pizca de honestidad de parte de los humanos. Esa bomba solo tiene un propósito: el exterminio de una especie. ¿Por qué construirían un arma así? Ilústreme, por favor.

Sentí la boca seca y escuché mi propia sangre martilleando en mis oídos. ¿Por qué la Federación tuvo que presenciar esa atrocidad? Los representantes compartían expresiones de horror y el pánico empezaba a cundir. Oí frases llamándonos monstruos y demonios.

Respiré profundo, buscando las palabras. —Lo construimos porque sabíamos que algún día el exterminio podría ser necesario. ¡Mire, no había otra forma!

—Quizás sea cierto, pero no siempre lo necesario es lo correcto —afirmó él.

—Entonces cree que nos equivocamos. Cree que somos monstruos como todos los demás, ¿no? —me levanté, con una ira ardiente recorriéndome. Todo este juicio era una trampa. No era de extrañar que Retke estuviera tan calmado; tenía esa grabación guardada todo el tiempo—. ¿Qué es lo que busca? ¿Que expulsen a la humanidad? ¿Que me encarcelen por crímenes de guerra?

—No. En realidad iba a pedirle que liderara nuestro ejército —Retke levantó una pata pidiendo silencio mientras estallaban las protestas—. Amigos, los humanos no pueden evitar lo que son, pero están intentando cambiar. Nunca nos han hecho daño. Lograron convencernos de que fueron la especie más pacífica de la galaxia por siglos.

—Todos hemos cometido errores. Quiero que avancemos juntos, sin que el odio nos detenga. Creo que este ser humano y sus compañeros nos harán más fuertes y seguros.

Miré a Retke con incredulidad. —No sé qué decir.

—Puede pensarlo, pero debe saber que si acepta... este patrón de mentiras y omisiones debe terminar aquí. Y por el bien del universo, no inventen más bombas. Creo que la humanidad ya tiene más que suficiente. Me gustaría que la galaxia siguiera existiendo mañana, ¿sabe?

Tal vez existiera un futuro donde volviéramos a ser amigos de la Federación. Las cosas serían distintas... pero quizás debían serlo. Tenía la esperanza sincera de no volver a usar nuestras peores armas.

Yo también quería que la galaxia existiera mañana.